\chapter{Presentación de análisis, resultados y discusión}

\section{Análisis descriptivo de la muestra}

[Presenta un análisis descriptivo detallado de la muestra utilizada en tu investigación.
Incluye características demográficas, estadísticas descriptivas y cualquier otra información
relevante que permita comprender las características de tu muestra.]

\section{Resultados obtenidos}

[Presenta los resultados principales de tu investigación. Orgánizalos de forma clara y lógica,
posiblemente mediante subsecciones. Utiliza tablas, figuras y gráficos para presentar los datos
de forma clara. Cada tabla o figura debe tener un título descriptivo y una referencia en el texto.]

\section{Discusión y análisis de los resultados}

[Analiza e interpreta los resultados obtenidos en el contexto del marco teórico y los objetivos
de investigación planteados. Discute:

\begin{itemize}
    \item Cómo se relacionan los resultados con investigaciones previas
    \item Las implicaciones teóricas de los hallazgos
    \item Las limitaciones de tu estudio
    \item Las posibles explicaciones de los resultados
    \item La significancia práctica de los hallazgos
\end{itemize}
]
